\chapter{Genomic sequence applications}
\label{evod48:chapter}
Two windows from the same locus can look like two examples while carrying
almost the same evidence. Suppose window A covers bases $[0,100)$ and B
covers $[80,180)$. A random split can train on A and test on B, sharing twenty
bases and possibly the very motif that defines the label. A larger model
cannot repair the resulting scientific question. We begin with the unit of
generalization, then make sequence orientation, metrics and interpretation
equally explicit.

DNA is biological material; a recorded DNA sequence is symbolic data; a
predictor operating on that sequence is a computational hypothesis. These
three statements do not collapse into one another. The executable work here
checks a genomic evaluation harness. It trains no genomic model and reports
no measured molecular effect.

\section{From a biological question to an observation contract}
\index{genomics!task contract}
Write the task as $(X,Y,U,M)$: permitted sequence and metadata $X$, target
$Y$, sampling unit $U$, and metric $M$. A promoter classifier predicts a
window label; a splice-site segmenter predicts positions; a chromatin-profile
model predicts assay tracks; a variant scorer compares two alleles in a
declared context. They have different losses, resolutions and information
boundaries. An offline chromatin predictor may use both flanks. A causal
next-base predictor may not inspect its target through an overlapping token.

For a promoter application, specify the organism, assembly, cell condition,
assay, positive-label rule and negative-sampling process. Promoters are DNA
regions at which transcription machinery associates to initiate transcription
\cite{evod48:promoter}. Sequence is only one part of that physical process.
The diagram distinguishes base pairing, protein association and transcription;
none of its arrows is an edge in a neural network.

\Plate{01-mechanism}{A schematic transcription setting supplies the biological
question. Solid paired rails denote DNA backbones, dashed rungs complementary
bases, and the oval associated transcription machinery. The RNA product grows
5$'$ to 3$'$. Sequence-based classification observes a window and an assay-derived
label; it does not simulate the machinery or establish its binding mechanism.}
{evod48:mechanism}

Each observation retains assembly, contig, zero-based half-open interval,
orientation, assay/source identifier, preprocessing version and label provenance.
The reference's \texttt{Example} carries the coordinate and grouping subset;
an actual data manifest must supply the remaining provenance. Its interval length
must equal its stored sequence length. Coordinates from two assemblies cannot
be compared by merely matching their integers. A liftover is a separately
versioned mapping with ambiguous and failed mappings reported.

\section{Leakage is a relation, then a transitive closure}
\index{data leakage!genomic overlap}\index{homology}\index{split closure}
The generalization question decides which relationships must remain inside a
split: overlap, locus, donor, homolog family, experiment or species. No split
is universally correct. An unseen-chromosome test asks a different question
from an unseen-person test on familiar loci. Both may still contain homologs
of training sequence. Pretraining exposure must also be audited when a
foundation model has already seen the reference genome.

Our small harness links two records when their same-assembly, same-contig
intervals overlap, or when their nonempty declared homology IDs agree.
Let $G=(V,E)$ be this undirected graph. Legal groups are its connected
components. For half-open intervals, overlap means
\[
\max(a_i,a_j)<\min(b_i,b_j).
\]
Adjacent $[0,4)$ and $[4,8)$ windows do not overlap. If A overlaps B and B
is homolog-linked to C, all three stay together even if A and C have no
direct link. A pairwise check on only the initially selected test IDs misses
this bridge.

\Plate{02-components}{Edges encode prohibited split crossings. Holding out A
also holds out B and C through transitive closure. D begins at C's end and has
no homology link, so adjacency alone does not join it. The grouping policy is
an input assumption, not a discovery of all possible biological dependence.}
{evod48:components}

\Listing{group_split}
\begin{proposition}[No declared edge crosses the split]
If each connected component is assigned wholly to one split, no edge in $E$
has endpoints in different splits.
\end{proposition}
\begin{proof}
The endpoints of an edge are connected by a path of length one and hence
belong to the same component. Whole-component assignment gives them the same
split. Transitive paths make the same argument apply to indirectly related
records.
\end{proof}
This theorem covers the declared relation only. Undetected paralogs,
misidentified donors or target-derived preprocessing can still leak. The
reference uses an $O(n^2)$ pair scan to expose semantics; production-scale
interval sweeps and homology clustering need their own parity tests.

\section{Reverse complement acts on data and labels}
\index{reverse complement}\index{unknown bases}\index{strand equivariance}
For $\Sigma=\{A,C,G,T,N\}$, complement exchanges A/T and C/G and fixes N.
For length $L$, $R(x)_i=c(x_{L-1-i})$. Applying reversal and complement twice
restores every index and base, so $R^2=I$. This is a representation symmetry,
not an assertion that all assays are orientation invariant.

For a scalar strand-invariant score, a shared two-pass construction is
\[
f_{\mathrm{sym}}(x)=\tfrac12[g(x)+g(Rx)].
\]
Then $f_{\mathrm{sym}}(Rx)=\tfrac12[g(Rx)+g(R^2x)]=f_{\mathrm{sym}}(x)$.
The equality assumes exactly the same deterministic $g$ and preprocessing in
both passes. A strand-specific label instead requires its declared label
action, potentially reversing positions and swapping strand channels. Pooling
it to one invariant scalar can erase the required distinction.

\Plate{03-strands}{Both strings describe the same antiparallel duplex, each read
5$'$ to 3$'$. A substitution at forward index $p$ moves to $L-1-p$ with complemented
alleles. Reversing positions without complementing bases produces the wrong
variant. Unknown N is an observation symbol, not a fifth chemical base.}
{evod48:strands}

The original package used CG as its motif. CG is its own reverse complement,
so that particular invariance test could pass without exercising two distinct
motif orientations. We keep the motif-count example but use ACG and its
reverse complement CGT. The assigned score for AACGTT is one, and the C-to-A
change at index two removes both contributions, giving $\Delta=-1$.
This is a detector calculation, not a learned affinity or expression change.

The reference rejects characters outside ACGTN, treats N as an all-zero
one-hot row with a separate validity flag, and returns an undefined GC fraction
for an all-N record. Silently treating N as A would fabricate a base. Chapter
23 accepts canonical ACGT; the local integration test uses that exact interface
and its offline shared-RC readout, with two untrained layers, learned absolute
positions and CPU float64. It proves neither learned invariance nor accuracy.

\section{Resolution, context and the baseline ladder}
\index{tokenization!base resolution}\index{long context}\index{baseline controls}
Base tokens preserve nucleotide resolution. Fixed $k$-mers change the sequence
length and perturbation footprint: one substituted base can affect several
overlapping tokens. Chapter 24's span availability condition remains necessary
even under triangular token attention. A context measured in tokens must also
be reported in bases, with stride, padding and uncovered tails declared.

Nucleotide Transformer evaluates pretrained Transformer representations on
genomic tasks; its original tokenization uses six-base units
\cite{evod48:nt}. HyenaDNA studies single-nucleotide modeling with contexts
up to one million tokens \cite{evod48:hyena}. Caduceus constructs
bidirectional RC-equivariant MambaDNA components \cite{evod48:caduceus}.
These are distinct prior approaches, not interchangeable instances of DOGMA.
The 2025 DNA-model benchmark evaluates frozen foundation representations
with downstream classifiers and finds task-dependent comparisons with
specialized models \cite{evod48:benchmark}. Frozen representations versus
task-trained competitors is an adaptation condition that must accompany
the conclusion; it is not a universal ranking of all fine-tuning procedures.

Start with composition/Markov or $k$-mer counts, linear and MLP readouts,
then local CNNs, GRUs, appropriate SSM/long-convolution models and Transformers.
Choose meaningful controls rather than requiring every architecture on every
task. Freeze splits, legal evidence, tuning trials, training tokens and primary
metric. Report parameter, state, context and measured serving budgets separately.
The actual Chapter 31 recurrence is a restricted teaching model, not full
DOGMA. Hermon DNA is Transformer-based; DOGMA is the non-Transformer direction.

A useful long-context ablation retains the central window while masking or
shuffling distal context under a declared nuisance-matching rule. Compare
several context lengths at matched training and tuning budgets. If only
same-locus random splits benefit, leakage is a competing explanation. If a
long context has no task information, success at storing it is not evidence
of regulatory understanding.

\section{Derive discrimination and probability checks}
\index{Matthews correlation coefficient}\index{average precision}\index{Brier score}
For binary predictions, retain $TP,TN,FP,FN$ before computing any scalar.
An all-negative classifier at 1\% prevalence has 99\% accuracy and zero recall.
The Matthews correlation coefficient is
\[
\mathrm{MCC}=\frac{TP\,TN-FP\,FN}
{\sqrt{(TP+FP)(TP+FN)(TN+FP)(TN+FN)}}.
\]
A zero denominator makes this expression undefined. Our harness reports
\texttt{None}; a benchmark choosing zero must disclose that convention.
The Brier score for probabilities $p_i$ is $n^{-1}\sum_i(p_i-y_i)^2$.
It measures squared probability error; a reliability curve by score bins
adds diagnostic information that one scalar cannot show.

Average precision (AP) is one specified summary of a precision--recall curve,
not an unspecified trapezoidal AUPRC. At distinct score thresholds $k$, let
$r_k$ be recall and $q_k$ precision after all tied examples enter. Define
\begin{equation}\label{evod48:ap}
\mathrm{AP}=\sum_k(r_k-r_{k-1})q_k.
\end{equation}
For scores $(.9,.8,.7)$ and labels $(1,0,1)$ the positive increments are
$1/2$ at precision one and $1/2$ at precision $2/3$: AP is $5/6$.
An all-tied predictor has AP equal to prevalence. Reordering tied rows cannot
improve it. With no positives AP is undefined under this contract.
\Listing{average_precision}
\begin{center}\input{\Generated/metrics.tex}\end{center}
The generated table uses labels $(1,1,0,0)$, predictions $(1,0,1,0)$ and
probabilities $(.8,.4,.6,.1)$. Accuracy is $1/2$, MCC zero, Brier $.1925$,
and AP $5/6$. These values validate arithmetic; they are not dataset results.

\section{Uncertainty follows the independent unit}
\index{cluster bootstrap}\index{sampling unit}
Holding trained predictors fixed, a cluster bootstrap resamples biological
groups with replacement and includes all windows of every sampled group.
If groups are independent draws from the target population, this approximates
sampling variation at that level. It does not include training-seed variation,
unseen assays or uncertainty about the grouping relation.
\Listing{cluster_interval}
The reported statistic is row-weighted accuracy on each resampled collection;
large groups receive more weight. Equal-group performance would first average
within groups and then across groups, a different estimand. Percentile endpoints
use the explicitly coded order-statistic indices. Two clusters are enough to
exercise the implementation and far too few for a convincing scientific
interval. Compare model differences on paired resamples, preserve seed-level
results, and stratify by distance, GC, ambiguous-base rate and label source.

\section{Variant scores and mechanistic claims}
\index{variant effect score}\index{biological validation}
For a normalized variant and fixed context, define
$\Delta=s(x_{\mathrm{alt}})-s(x_{\mathrm{ref}})$.
The sign convention, score scale, reference allele, assembly and context must
be fixed before interpretation. The local oracle implements substitutions only;
indels need an explicit alignment and output-coordinate policy.
\Listing{snv}
Saliency, attention and in-silico mutation measure sensitivity of a predictor.
Controls include label shuffling, matched GC negatives, motif insertion/deletion,
RC consistency and perturbations outside the tested footprint. A motif effect
in a synthetic detector is a useful positive control for the harness. It does
not show that a protein binds that sequence in a cell.

A biological follow-up protocol states the hypothesized mechanism, intervention,
assay, cell condition, replication unit and alternative explanations. Reference
and alternate constructs need matched backgrounds, assay controls and an analysis
plan. No wet-lab experiment is conducted here. Computational prioritization is a
handoff to such evidence, not a replacement for it.

\section{Selective glossary}
\begin{description}
\item[Sampling unit.] The biological entity treated as independent for the
generalization and uncertainty claim; it need not equal one sequence window.
\item[Split closure.] Connected-component grouping under all declared relations
that must not cross train/evaluation boundaries.
\item[RC action.] Reversal plus base complement, accompanied by the appropriate
coordinate and label transformation.
\item[Average precision.] Recall-increment-weighted precision with a declared
tie policy; different from an unspecified trapezoidal PR area.
\item[Variant score.] A model-score contrast between normalized reference and
alternate sequence, with no automatic causal or clinical interpretation.
\end{description}

\section{Exercises and worked solutions}
\begin{enumerate}
\item Why does the opening random split change the question? \textbf{Answer:}
it tests a partly observed locus rather than an independently unseen locus.
\item Do adjacent half-open intervals overlap? \textbf{Answer:} no; their
intersection has zero length and the strict inequality fails.
\item A overlaps B; B and C share a family. Hold out A. \textbf{Answer:}
hold out A, B and C by connected-component closure.
\item Does the closure theorem prove all homology is removed? \textbf{Answer:}
only links supplied by the audited relation; undetected homology remains possible.
\item Compute $R(\mathrm{AACGTN})$. \textbf{Answer:} NACGTT.
\item Prove $R^2=I$. \textbf{Answer:} each index reverses twice and each
base is complemented twice.
\item Why is CG a weak two-orientation test motif? \textbf{Answer:} its RC
is CG, so missing the second orientation can be hidden.
\item Transform a substitution at $p=2$ in a six-base record. \textbf{Answer:}
index three, with both reference and alternate complemented.
\item Why keep a validity mask for N? \textbf{Answer:} zero feature mass must
not be confused with observed adenine or an actual chemical fifth base.
\item Must a splice track be invariant? \textbf{Answer:} no; positional and
strand channels can require an equivariant label transformation.
\item A 1\% positive task, all-negative predictor? \textbf{Answer:}
99\% accuracy, zero recall, undefined MCC under our denominator convention.
\item Compute AP for $(1,0,1)$ sorted by descending score. \textbf{Answer:}
$\tfrac12+\tfrac12(2/3)=5/6$.
\item What is all-tied AP at 20\% prevalence? \textbf{Answer:} $.2$;
breaking ties by labels would leak evaluation information.
\item Brier for labels $(1,0)$ and probabilities $(.8,.2)$? \textbf{Answer:}
$(.04+.04)/2=.04$.
\item Does resampling windows estimate donor uncertainty? \textbf{Answer:}
not when windows from one donor are dependent; resample the donor groups.
\item Why distinguish frozen embeddings and fine-tuning? \textbf{Answer:}
they permit different adaptation and compute, changing the comparison.
\item Is the tested Hermon path a newly trained genomics model? \textbf{Answer:}
no; it is the actual untrained Chapter 23 interface checked for algebraic symmetry.
\item Design a distal-context ablation. \textbf{Rubric:} central evidence fixed,
declared distal perturbation, nuisance matching, legal split, matched opportunity,
paired metrics, seed/cluster uncertainty and retained negative results.
\item What could falsify the motif explanation? \textbf{Answer:} gain survives
motif deletion but vanishes with a GC-matched nuisance intervention, motivating
a competing shortcut explanation.
\item Propose a biological validation handoff. \textbf{Rubric:} explicit allele
and cell context, causal intervention, assay controls, independent replication,
predeclared analysis and bounded interpretation; no invented measurements.
\end{enumerate}

\section{Bridge}
Chapter 49 changes domains but preserves the central discipline: record what
was observed, retain the identity of carried state, and declare which evidence
may influence the next decision. A persistent service needs those contracts
even when its inputs are language and code rather than genomic bases.
